\section{Implementación}

\subsection{Diseño del sistema}
UrbanTracker está dividido en frontend y backend. En el lado de usuario hay tres piezas:

Web pública para consultar rutas y ver los buses en el mapa.

App móvil del conductor (React Native) que toma el GPS del teléfono.

Panel web para administración de rutas, vehículos y conductores.

El backend está hecho en Spring Boot y expone servicios REST. La API está documentada con OpenAPI/Swagger para facilitar pruebas e integración \cite{easyrest2022}.

\subsection{Arquitectura del backend (hexagonal y modular)}
El backend sigue una arquitectura hexagonal organizada por módulos.

domain: reglas del negocio (rutas, vehículos, conductores, recorridos).

application: funciones que se encargan de cada acción del sistema (por ejemplo, iniciar un recorrido o procesar una ubicación recibida).

infrastructure: adaptadores que publican REST, hablan con PostgreSQL/JPA y gestionan la comunicación en tiempo real.

Esta separación ayuda a mantener el código claro y fácil de extender (POO y separación de responsabilidades) \cite{poo2023}.

\subsection{Aplicaciones de usuario}
Web pública (React/Next.js): muestra rutas y los buses moviéndose en el mapa, y permite filtros básicos.

Panel administrativo (React/Next.js): alta/edición de rutas, vehículos y conductores, con una vista general de la flota.

App del conductor (React Native): pide permisos de ubicación, permite iniciar/cerrar recorrido y envía periódicamente su ubicación.

\subsection{Comunicación en tiempo real}
El tiempo real funciona con un flujo simple:

El teléfono del conductor envía su posición por MQTT a un broker.

El backend recibe esos mensajes, los revisa y los reenvía por WebSockets a los clientes de la web \cite{sismo2016}.

Este enfoque reduce esperas para el usuario y evita depender de consultas periódicas. La utilidad de mostrar posiciones en vivo está documentada en trabajos previos \cite{cityruta2018}; \cite{stopbus2016}.

\subsection{Base de datos}
Se usa PostgreSQL para las entidades operativas del sistema (rutas, vehículos, conductores, asignaciones/recorridos). No se guardan trazas de posición en esta etapa.

\subsection{Seguridad y control de acceso}
El acceso a operaciones de conductor y administración se protege con JWT. Los roles limitan lo que cada perfil puede hacer desde la web o el panel.

\subsection{Integración de mapas}
Para la visualización se integra Mapbox en modo 3D en la web. Esto ayuda a entender mejor el movimiento de cada bus sobre su ruta y mejora la experiencia de quien consulta.